\documentclass[12pt]{article}
\usepackage[T1]{fontenc}
\usepackage[utf8]{inputenc}
\usepackage{lmodern}
\usepackage{textcomp}
\usepackage{enumitem}
\usepackage{geometry}
\usepackage{graphicx}
\usepackage{amsmath, amssymb}
\geometry{margin=1in}
\begin{document}
\begin{center}
Physics - Graduate Level Assessment 20 \
Total Marks: 40 \
Answer all questions.
\end{center}

\section*{Multiple Choice (10 marks)}
\begin{enumerate}[label=\arabic*.]
\item The half-life of radon is 3.80 days. After how many days will only one-sixteenth of a radon sample remain?
\begin{enumerate}[label=(\alph*)]
    \item 26.6 days
    \item 19.0 days
    \item 3.80 days
    \item 5.7 days
    \item 22.8 days
\end{enumerate}
\item What is the intensity of the electric field 10 cm from a negative point charge of 500 stat-coulombs in air?
\begin{enumerate}[label=(\alph*)]
    \item 5 dyne/stat-coul
    \item 8 dyne/stat-coul
    \item 0.5 dyne/stat-coul
    \item 100 dyne/stat-coul
    \item 10 dyne/stat-coul
\end{enumerate}
\item Calculate the mean ionic activity of a $0.0350 \mathrm{~m} \mathrm{Na}_3 \mathrm{PO}_4$ solution for which the mean activity coefficient is 0.685 .
\begin{enumerate}[label=(\alph*)]
    \item 0.0612
    \item 0.0426
    \item 0.0753
    \item 0.0905
    \item 0.0294
\end{enumerate}
\item determine the ratio of the radius of a uranium-238 nucleus to the radius of a helium-4 nucleus.
\begin{enumerate}[label=(\alph*)]
    \item 4.8
    \item 7.1
    \item 5.5
    \item 1.8
    \item 4.2
\end{enumerate}
\item The Pleiades is an open star cluster that plays a role in many ancient stories and is well-known for containing ... bright stars.
\begin{enumerate}[label=(\alph*)]
    \item 5
    \item 7
    \item 9
    \item 12
\end{enumerate}
\end{enumerate}

\section*{True / False (10 marks)}
\textit{Answer True or False}
\begin{enumerate}[label=\arabic*.]
\setcounter{enumi}{5}
\item True or False: The correct answer to 'What equal positive charges would have to be placed on Earth and on the Moon to neutralize their gravitational attraction?' is '$3.3 \times 10^{13} \mathrm{C}$'.
\item True or False: Meteorites always exhibit a high density due to their unique composition and formation process.
\item True or False: The potential energy of a compressed spring and a charged object depends solely on the distance they travel.
\item True or False: A completely submerged object displaces a volume of fluid equal to the pressure exerted by the fluid on its surface.
\item True or False: A voltage of 10 volts is sufficient to send a current of 2 amperes through a circuit with a resistance of 3 ohms.
\end{enumerate}

\section*{Long Form Response (20 marks)}
\textit{Answer each question with a clear and complete explanation.}
\begin{enumerate}[label=\arabic*.]
\setcounter{enumi}{10}
\item Calculate the net electric flux through a Gaussian cube with an edge length of 55 cm, centered around a particle with a charge of 1.8 µC. Discuss the principles involved.
\item A child slides a block of mass 2 kg at an initial speed of 4 m/s towards a spring with a spring constant of 6 N/m. Analyze the energy transformations involved and determine the maximum compression of the spring, providing a detailed explanation of your reasoning and calculations.
\end{enumerate}

\vfill
\noindent\textit{End of Paper}
\end{document}